% !TEX root = ../thesis.tex

MicroRNAs (miRNAs) are pivotal post-transcriptional regulators, modulating up to 60\% of protein-coding genes to maintain cellular homeostasis and influence disease pathogenesis. Despite their centrality, miRNA research remains siloed: upstream transcriptional annotations are incomplete, network analyses prioritize node-centric expression over edge relationships, and oxidative epitranscriptomic modifications are largely overlooked. These limitations lead to poor reproducibility and impede the translation of miRNA findings into robust clinical biomarkers. To overcome these challenges, this thesis constructs a systematic framework that transitions miRNA studies from fragmented, static, molecule-focused approaches to a holistic, dynamic, network-oriented paradigm, enabling mechanistic investigation and predictive modeling based on miRNA-mediated regulatory networks during disease progression.

First, addressing the upstream gap, we developed miRStart 2.0, a deep learning-based platform integrating over 4,500 multi-modal datasets. Achieving high-resolution accuracy against experimentally validated data, it annotated 28,828 miRNA transcription start sites and mapped over 6 million tissue-specific TF-miRNA interactions, laying a comprehensive data foundation for transcriptional control. Second, elucidating network architecture, we demonstrated that ``edge information'' (regulatory interactions) offers insights equal to node abundance. In triple-negative breast cancer, our interaction-based prognostic model significantly outperformed traditional expression-based methods. Furthermore, topological analysis of the hub miR-483 resolved its context-dependent functional duality across diverse diseases, proving that network wiring dictates function. Third, unveiling oxidative rewiring, we developed a computational pipeline to detect oxidative footprints from standard miRNA-seq data. We introduced the rewiring activity score to quantify this dynamic network reprogramming, identifying stage-specific biomarkers in liver cancer that drive a critical metabolic shift toward oxidative phosphorylation.

In conclusion, this framework integrates transcriptional origins, topological architecture, and epitranscriptomic plasticity, resolving the isolation in current miRNA research. By capturing the full complexity of miRNA-mediated regulatory networks, this work advances a systems-level comprehension of gene regulation and paves the way for novel precision interventions in complex diseases.
